\section{Discussion}\label{sec:discussion}

\subsection{Related work}\label{sec:related}

Approaches to a location degeneracy remove the directions (constraints,
projections, integration), reparameterize them (centering, non-centering,
interweaving), or adapt the sampler's metric to them. The rotation is a
reparameterization chosen so that the third becomes cheap, and unlike
removal it keeps the original parameter space and posterior and returns
the original parameters.

\paragraph{Choosing a parameterization.} Centering
\citep{gelfand1995,gelfand1996} removes the location degeneracy and
non-centering the funnel \citep{papaspiliopoulos2003,papaspiliopoulos2007};
partial non-centering interpolates, the ancillarity-sufficiency
interweaving strategy \citep{yu2011} interweaves the
two within a Gibbs sampler, and \citet{gorinova2020} automate the choice
per variable. All of these choose between centered and non-centered
coordinates for each random effect. The rotation
keeps the non-centered parameterization, with its benefits for the
variances, and changes only which linear combinations of the standardized
effects are sampling coordinates; against JMbayes2, which centers the
coefficients it can, it achieved comparable efficiency on those
coefficients (Section~\ref{sec:realdata}). \citet{browne2009} studied
reparameterizations for multilevel survival models under Gibbs sampling.
Their orthogonal parameterization orthogonalizes the fixed-effect
predictors among themselves, as Stan's QR reparameterization does
\citep{stanug}, but they observe that predictors correlated with the
cluster indicator vectors make fixed and random effects ``play similar
roles'' and suggest orthogonality to those indicators as future work. The
absorbable subspace is that span, handled by rotating the random effects
rather than the predictors.

\paragraph{Removing the directions.} \citet{zanella2021} show that
sum-to-zero constraints always improve the Gibbs sampler for crossed
random effects, and that they also help NUTS dramatically on a Poisson
crossed model (minimum ESS from 3.7 to 4977). Their construction samples
the posterior conditioned on the constraint, and Stan's
\texttt{sum\_to\_zero\_vector} likewise changes the prior unless
compensated. Restricted spatial regression \citep{reich2006,hughes2013}
projects spatial random effects off the fixed-effect design, and is known
to affect inference on the fixed effects
\citep{hodges2010,hanks2015,khan2022}. An earlier version of
\texttt{jmjax} swept the absorbable directions out of the sampled
random-effect coordinates: without a correction this miscalibrated the
intervals for the swept coefficients, and with an exact correction it was
about as efficient as the rotation in ordinary designs and up to twice as
efficient in the most information-dense ones (Section~\ref{sec:stress}),
at the price of the correction step.
The closest precedent is \citet{vines1996}, whose transformation is exact:
for independent Gaussian random effects and a Gaussian prior on the grand
mean, the directions that shift the grand mean against the mean of each
factor appear in no likelihood term and can be integrated out. Their
method covers the grand-mean direction only, reports sum-to-zero contrasts
rather than the original parameters, and leaves correlated random effects
to further work; they also warn that recentering the effects inside each
iteration has an unknown effect on the stationary distribution. The
rotation extends the exact approach to every absorbable direction,
including covariate-weighted and random-slope ones, to correlated random
effects and to any prior on the fixed effects.

\paragraph{Analyzing the directions.} In the discussion of
\citet{zanella2021}, \citet{yang2021} decompose the Gibbs sampler for a
random-intercept model with covariates using the subject-level covariate
means, which span the covariate-weighted absorbable subspace for the
intercept column, and show that the convergence rate is governed by the
correlation between $\beta$ and the random effects along it; they note
that the decomposition breaks down for general designs, and do not treat
random slopes. In the same discussion \citet{robert2021} asked whether
these parameterization choices had been studied for HMC, and the
rejoinder \citep{zanella2021} called that open. Lemma~\ref{lem:local} and
Proposition~\ref{prop:ridge} are an analysis on the HMC side for the
absorbable directions. They complement general results on linear
preconditioning for Markov chain Monte Carlo \citep{hirdlivingstone2025}: the block the
rotation creates is small enough to precondition densely because the model
says where the ridge is.

\paragraph{Joint-model software.} JMbayes2 \citep{jmbayes2} uses
Metropolis-within-Gibbs; \texttt{stan\_jm} in rstanarm
\citep{brilleman2018} uses NUTS with non-centered random effects and so,
we expect, meets the same ridge, which we have not tested; INLAjoint
\citep{rustand2024} uses integrated nested Laplace approximations. We have
not found a previous treatment that applies the data-augmentation
autocorrelation identity \citep{liu1994} to the association parameter
together with a measurement of how its shared variance divides between the
random effects and the baseline hazard.

\subsection{Limits and extensions}

With about 9 visits per subject the rotation gave $\alpha$ 0.72--1.00 of
the unrotated fit's ESS per second, though it gained with denser follow-up;
the mechanism is not established. In the most information-dense designs
the swept variant was up to twice as efficient, as the fixed-metric
residual of Section~\ref{sec:residual} predicts, although the rotation was
still 5--30 times ahead of the unrotated fit there. The simulation
evidence comes from one design family with a spline baseline hazard. The
basis construction's rank decision is a tolerance judgment when
absorbable and identified directions are nearly collinear, although every
returned direction is verified before use, and Condition~\ref{cond:S} can
fail for very flexible time bases with few observations per subject.

A metric that follows the sampled $\sigma_0$, for example by rescaling the
dense block with the current $L$, is the obvious way to remove the
residual and has not been tried. The construction needs only the fixed-
and random-effects designs, so it applies to any model whose random
effects enter through a linear predictor, including the GLMMs of brms and
rstanarm. Crossed random effects need one basis per grouping factor, each
rotating that factor's own index; the argument of
Proposition~\ref{prop:exact} carries over unchanged.

\subsection{Conclusion}

The location degeneracy between fixed and random effects is shared by
every regression with a random intercept and subject-constant fixed-effect
columns, and it recurs for the time slope in every model with a random
slope; this much is algebra. In the models studied here it slowed HMC
substantially on the coefficients an analyst reports. Rotating the
standardized random effects so that the absorbable directions become a
small block of coordinates, and giving that block and the fixed effects
one dense metric block, largely whitens the resulting ridge without
changing the model. The association parameter of a joint model is a
separate matter: in the two joint models measured, its posterior variance
was largely shared with the survival parameters, a sharing that a
sampler updating it conditionally pays for in autocorrelation, whereas HMC,
heuristically, pays for it, if at all, in integration steps.
